\chapter*{Recomendaciones}
\addcontentsline{toc}{chapter}{Recomendaciones}\label{chap:recommendations}

A partir de este trabajo, y tras considerar tanto las limitaciones que la sección anterior enumera como las oportunidades de extensión del sistema, el capítulo formula las recomendaciones siguientes para trabajos futuros.

\begin{enumerate}[leftmargin=*,labelsep=0.5em]
	\item \textbf{Proveedores de autenticación delegada.} Integrar mecanismos de autenticación delegada (OAuth con Google o Microsoft, integración con LDAP institucional) en la aplicación web, con vistas al escenario en que el sistema opere dentro de la red institucional. La integración facilitaría el alta de usuarios desde el directorio institucional, eliminaría la duplicación de credenciales y permitiría aprovechar las políticas de acceso que la organización ya aplica.
	
	\item \textbf{Infraestructura para modelos múltiples.} Implementar una infraestructura de soporte para varios modelos OCR que permita registrar nuevos modelos sin alterar el resto del código y que combine las predicciones de los modelos disponibles —mediante votación, promediado de probabilidades a nivel de carácter o decodificación conjunta— para mejorar la robustez frente a variaciones tipográficas y de degradación. La infraestructura facilitaría además la sustitución del modelo activo sin interrupción del servicio y la comparación directa entre versiones sucesivas del modelo.
	
	\item \textbf{Cola de revisiones pendientes.} Incorporar una cola de revisiones que permita asignar trabajo de transcripción a un usuario concreto, a un subconjunto de usuarios o al colectivo entero. La cola debe acoplar al resto de operaciones del sistema —eliminación de documentos, cambios de versión, asignación de roles— para mantener la coherencia del estado de la base de datos y evitar referencias huérfanas tras los borrados, y debe ofrecer una vista por usuario que sintetice la carga pendiente y el historial de revisiones que el usuario ya terminó.
	
	\item \textbf{Aprendizaje continuo a partir de las transcripciones revisadas.} Establecer un ciclo de aprendizaje continuo en el que las transcripciones que los revisores corrigen sobre la aplicación pasen a integrar el \textit{corpus} de entrenamiento, y planificar etapas periódicas de adaptación del modelo sobre el \textit{corpus} que las revisiones vayan acumulando. El ciclo permitiría especializar el modelo en las tipografías, los registros léxicos y los patrones de degradación característicos del acervo de la institución, y reduciría el coste manual de revisión a medida que el modelo mejora sobre los documentos propios.
	
	\item \textbf{Cola persistente de procesamiento por lotes.} Migrar el procesamiento OCR a una cola persistente que descanse sobre Celery~\cite{celery}, con Redis~\cite{redis} o RabbitMQ~\cite{rabbitmq} como gestor de mensajes. La cola permitiría procesar grandes volúmenes de páginas sin bloquear la interfaz web, soportar la cancelación y la reanudación de trabajos, y aprovechar varios trabajadores en paralelo cuando el volumen lo justifique. La persistencia del estado de la cola sobreviviría además a reinicios del servidor sin pérdida de trabajos en curso.
\end{enumerate}

Las contribuciones del trabajo —el \textit{corpus} sintético, el \textit{pipeline} de preprocesamiento configurable y la infraestructura web de revisión— pueden ayudar a otros proyectos de digitalización documental que enfrenten el mismo conjunto de problemas: escasez de datos para el idioma de trabajo, diversidad tipográfica y necesidad de revisión colaborativa. La extensión del sistema a \textit{corpus} en otras lenguas romances (catalán, portugués, gallego) requeriría únicamente la sustitución del \textit{corpus} textual base y de la distribución de frecuencias de caracteres objetivo, sin cambios en el \textit{pipeline} ni en la arquitectura del modelo. La extensión a otros tipos documentales —partituras musicales, fórmulas químicas, dibujos técnicos— exigiría rediseñar el vocabulario y la simulación de degradación, pero el resto del \textit{pipeline} permanecería operativo.

La incorporación del sistema a la práctica institucional del archivo del Instituto Cubano de Investigación Cultural Juan Marinello requiere tres condiciones: disponibilidad de un servidor con capacidad de cómputo adecuada para la inferencia OCR en lotes, capacitación del personal del archivo en el flujo de revisión y validación de transcripciones, y un protocolo de control de calidad que combine la inspección humana con las métricas automáticas del sistema. Una vez que la institución cumpla las tres condiciones, el sistema permitirá al archivo procesar volúmenes de documentación que el flujo manual actual no es capaz de absorber, y abrir el acervo a la consulta computacional, a la indexación y al análisis a escala.
